\documentclass[10pt,a4paper]{article}
\usepackage[T1]{fontenc}
\usepackage[utf8]{inputenc}
\usepackage{lmodern,microtype,amsmath,booktabs,graphicx,siunitx}
\usepackage[margin=20mm]{geometry}
\usepackage[colorlinks=true,urlcolor=blue,citecolor=blue]{hyperref}
\setlength{\parindent}{0pt}\setlength{\parskip}{5pt}
\setlength{\emergencystretch}{2em}
% Generated from analysis.h5; do not edit.
\newcommand{\IonWorkEv}{21.9}
\newcommand{\IonPitchMm}{2}
\newcommand{\IonFrameMs}{20}
\newcommand{\IonZBinMm}{2}
\newcommand{\IonSpeedMmMs}{0.1}
\newcommand{\IonMarginFraction}{0.1}
\newcommand{\IonTolerance}{1e-11}

\title{Initial ionization from NEXUS tracks}
\author{NAUSICAA analysis}\date{}
\begin{document}\maketitle
\section{Calculation and interpretation}
For each deposited-energy hit, the expected number of positive ions is
\[
 \overline n_i=10^6 E_i[\mathrm{MeV}]/W[\mathrm{eV}],\qquad
 E_{\rm evt}=\sum_i E_i,\qquad \overline N_{\rm evt}=\sum_i\overline n_i.
\]
We adopt $W=\IonWorkEv$ eV as a configurable nominal mean energy per electron--ion
pair in gaseous xenon. It describes the partition of deposited energy between
ionization and other processes, rather than the atomic ionization threshold;
see Ref.~\cite{Serra:2014zda}. This is an adopted input, not a measurement from these samples.
Counts are continuous expectations, not rounded integer realizations. Each
molecular $\mathrm{Xe}_2^+$ ion carries charge $+e$.

The per-event mean ions per hit is $\overline N_{\rm evt}/N_{\rm hits}$.
The table reports the mean and population dispersion of this quantity across
events. The pooled value instead divides total ions by total hits. Hit-level
quantities depend on the Geant4 deposition subdivision.

Each event is translated to the common generation vertex of its primary
particles, copied from the original NEXUS particle table:
$\mathbf r'_i=\mathbf r_i-\mathbf r_{\rm vertex}$. Orientations are preserved;
no energy barycenter is used to define the origin.
XY bins have pitch \IonPitchMm~mm. With the configured uniform speed
\IonSpeedMmMs~mm/ms, a \IonFrameMs~ms interval corresponds to a Z bin of
\IonZBinMm~mm. The time axis is relative to the generation vertex, not an absolute
arrival-time prediction. Projections are averaged over all saved events,
including zero contributions, without normalizing individual event energies.

This stage describes initial ion production. It includes no ionization
fluctuations, recombination, diffusion, focusing or collection losses.
The two marginal projections do not determine joint pad-and-frame occupancy.

\par The spatial selection keeps $|x'|\leq 150.0$ mm and the same bound on $|y'|$. It is applied after translating to the generation vertex. No Z cut or subsequent recentering is applied. Raw energy totals include every original hit; the mean maps use the selected hits and the original number of events. The full uncut marginal histograms and per-event discarded energy are retained. Common grids enclose every selected hit with the configured margin; bins are left-closed/right-open, with the final upper edge included.
Common XY grid: 168 by 168 bins, from -168 to 168 mm on each axis. Z: 1896 bins, from -1896 to 1896 mm.
\clearpage
\section{Double beta}
% Generated from analysis.h5
\newcommand{\IonSampleAEvents}{8995}
\newcommand{\IonSampleAHits}{2833001}
\newcommand{\IonSampleATotalIons}{1008874952.97}
\newcommand{\IonSampleAMeanIons}{112159.527846}
\newcommand{\IonSampleAMeanEnergyMeV}{2.45629365982}
\newcommand{\IonSampleAPooledIonsPerHit}{356.115283042}

Saved events: \num{8995}; hits: \num{2.833e+06}.
Total expected ions: \num{1.00887e+09}. Pooled ions/hit: \num{356.115}.
\begin{center}\small\begin{tabular}{lrrrr}\toprule
Quantity & Mean & Population SD & Minimum & Maximum \\\midrule
Energy (MeV) & \num{2.45629} & \num{0.0132068} & \num{2.30009} & \num{2.45783} \\
Expected ions/event & \num{112160} & \num{603.049} & \num{105027} & \num{112230} \\
Mean ions/hit in an event & \num{360.986} & \num{41.3303} & \num{211.754} & \num{505.539} \\
Hits/event & \num{314.953} & \num{37.6124} & \num{222} & \num{530} \\
\bottomrule\end{tabular}
\end{center}
The table above uses all original hits. After the vertex-relative XY cut, the mean retained energy is \num{2.41807} MeV and the mean retained ion count is \num{110414} per event.
Discarded hits: \num{84533} (\num{2.98387}\%). Discarded ions: \num{1.56993e+07} (\num{1.55612}\%). Empty events: \num{0}.
Maximum discarded ion fraction in an individual event: \num{70.9945}\%. Discarded energy is never restored by rescaling.
\begin{figure}[htbp]\centering
\includegraphics[width=.49\linewidth]{\detokenize{../../results/ion_tracks/0nubb/xy_mean.png}}
\includegraphics[width=.49\linewidth]{\detokenize{../../results/ion_tracks/0nubb/z_mean.png}}
\caption{Selected ensemble mean, with zero contributions included in the denominator. XY is integrated over z; Z is integrated over x and y. The right panel shows the central longitudinal detail; it applies no additional Z cut.}
\end{figure}
\begin{figure}[htbp]\centering
\includegraphics[width=.49\linewidth]{\detokenize{../../results/ion_tracks/0nubb/xy_event_0.png}}
\includegraphics[width=.49\linewidth]{\detokenize{../../results/ion_tracks/0nubb/z_event_0.png}}
\caption{Selected hits of example event 0. XY is integrated over z; Z is integrated over x and y. The right panel shows the central longitudinal detail; it applies no additional Z cut.}
\end{figure}
\clearpage\subsection*{Tail diagnostics}
\begin{figure}[htbp]\centering
\includegraphics[width=.82\linewidth]{\detokenize{../../results/ion_tracks/0nubb/xy_tails_hits.png}}
\caption{Uncut vertex-relative distributions. Dashed lines mark the XY cut; bottom panels show a detail of the same histograms.}
\end{figure}
\begin{figure}[htbp]\centering
\includegraphics[width=.82\linewidth]{\detokenize{../../results/ion_tracks/0nubb/xy_tails_ions.png}}
\caption{Uncut vertex-relative distributions. Dashed lines mark the XY cut; bottom panels show a detail of the same histograms.}
\end{figure}
\begin{figure}[htbp]\centering
\includegraphics[width=.82\linewidth]{\detokenize{../../results/ion_tracks/0nubb/z_mean_full.png}}
\caption{Full longitudinal distribution of retained hits.}
\end{figure}
\clearpage
\section{Single electron}
% Generated from analysis.h5
\newcommand{\IonSampleBEvents}{8483}
\newcommand{\IonSampleBHits}{2732400}
\newcommand{\IonSampleBTotalIons}{951436606.213}
\newcommand{\IonSampleBMeanIons}{112158.034447}
\newcommand{\IonSampleBMeanEnergyMeV}{2.45626095439}
\newcommand{\IonSampleBPooledIonsPerHit}{348.205462675}

Saved events: \num{8483}; hits: \num{2.7324e+06}.
Total expected ions: \num{9.51437e+08}. Pooled ions/hit: \num{348.205}.
\begin{center}\small\begin{tabular}{lrrrr}\toprule
Quantity & Mean & Population SD & Minimum & Maximum \\\midrule
Energy (MeV) & \num{2.45626} & \num{0.0127885} & \num{2.30124} & \num{2.45783} \\
Expected ions/event & \num{112158} & \num{583.948} & \num{105079} & \num{112230} \\
Mean ions/hit in an event & \num{353.119} & \num{41.0734} & \num{211.355} & \num{503.272} \\
Hits/event & \num{322.103} & \num{39.048} & \num{223} & \num{531} \\
\bottomrule\end{tabular}
\end{center}
The table above uses all original hits. After the vertex-relative XY cut, the mean retained energy is \num{2.39425} MeV and the mean retained ion count is \num{109326} per event.
Discarded hits: \num{121374} (\num{4.44203}\%). Discarded ions: \num{2.40203e+07} (\num{2.52463}\%). Empty events: \num{0}.
Maximum discarded ion fraction in an individual event: \num{89.5198}\%. Discarded energy is never restored by rescaling.
\begin{figure}[htbp]\centering
\includegraphics[width=.49\linewidth]{\detokenize{../../results/ion_tracks/electron_0nubb_energy/xy_mean.png}}
\includegraphics[width=.49\linewidth]{\detokenize{../../results/ion_tracks/electron_0nubb_energy/z_mean.png}}
\caption{Selected ensemble mean, with zero contributions included in the denominator. XY is integrated over z; Z is integrated over x and y. The right panel shows the central longitudinal detail; it applies no additional Z cut.}
\end{figure}
\begin{figure}[htbp]\centering
\includegraphics[width=.49\linewidth]{\detokenize{../../results/ion_tracks/electron_0nubb_energy/xy_event_0.png}}
\includegraphics[width=.49\linewidth]{\detokenize{../../results/ion_tracks/electron_0nubb_energy/z_event_0.png}}
\caption{Selected hits of example event 0. XY is integrated over z; Z is integrated over x and y. The right panel shows the central longitudinal detail; it applies no additional Z cut.}
\end{figure}
\clearpage\subsection*{Tail diagnostics}
\begin{figure}[htbp]\centering
\includegraphics[width=.82\linewidth]{\detokenize{../../results/ion_tracks/electron_0nubb_energy/xy_tails_hits.png}}
\caption{Uncut vertex-relative distributions. Dashed lines mark the XY cut; bottom panels show a detail of the same histograms.}
\end{figure}
\begin{figure}[htbp]\centering
\includegraphics[width=.82\linewidth]{\detokenize{../../results/ion_tracks/electron_0nubb_energy/xy_tails_ions.png}}
\caption{Uncut vertex-relative distributions. Dashed lines mark the XY cut; bottom panels show a detail of the same histograms.}
\end{figure}
\begin{figure}[htbp]\centering
\includegraphics[width=.82\linewidth]{\detokenize{../../results/ion_tracks/electron_0nubb_energy/z_mean_full.png}}
\caption{Full longitudinal distribution of retained hits.}
\end{figure}
\clearpage\section{Validation and provenance}
The Julia calculation checks that each event's XY and Z projections sum to its
retained expected ion count. The summed mean projections must equal the retained mean ion count
per event. No renormalization is applied. Out-of-range hits cause an error.
The following relative errors must be below \num{\IonTolerance}.
\begin{center}\begin{tabular}{lrrr}\toprule
Sample & Worst event & Mean XY & Mean Z\\\midrule
Double beta & \num{8.88e-16} & \num{2.22e-16} & \num{6.66e-16} \\
Single electron & \num{8.88e-16} & \num{1.11e-16} & \num{1.11e-16} \\\bottomrule\end{tabular}\end{center}
\subsection*{Double beta}
\begin{tabular}{lp{.65\linewidth}}\toprule
Input & events/0nubb\_hitdf.h5 \\
Original & 0nubb.next.h5 \\
Generator & Decay0Interface \\
Generated events & 10000 \\
Source lower energy cut & 2300 keV \\
Source upper energy cut & 2470 keV \\
Source radius & 2. m \\
Source height & 4. m \\
Source pressure & 15. bar \\
\bottomrule\end{tabular}
\par\small Input SHA-256: \texttt{627a22476784b5279e8d716b581237c50789ba07ba1c80d0a573abfcb8c6e12f}.
\par Analysis code SHA-256: \texttt{6ed06e3d0e26540893703c8659bba7f37a569929b1cceadec1b66ad86e2b2573}.
ormalsize
\subsection*{Single electron}
\begin{tabular}{lp{.65\linewidth}}\toprule
Input & events/electron\_0nubb\_energy\_hitdf.h5 \\
Original & electron\_0nubb\_energy.next.h5 \\
Generator & SingleParticleGenerator \\
Generated events & 10000 \\
Source lower energy cut & 2300 keV \\
Source upper energy cut & 2470 keV \\
Source radius & 2. m \\
Source height & 4. m \\
Source pressure & 15. bar \\
\bottomrule\end{tabular}
\par\small Input SHA-256: \texttt{8d36747bba5b67515a9bf19d3d94aeb7252b3bf77f2a6583a1564c8a6852ecbe}.
\par Analysis code SHA-256: \texttt{6ed06e3d0e26540893703c8659bba7f37a569929b1cceadec1b66ad86e2b2573}.
ormalsize
\par Numerical quantities, plots and tables are read from the Julia HDF5 result.
The HDF5 files also retain configuration and software hashes, per-event
quantities, bin edges and per-hit ion expectations in source row order.
\bibliographystyle{unsrt}\bibliography{references}
\end{document}
